\subsection{Composite Residuosity}\label{sec:trapdoor_functions:composite_residuosity}

The \textbf{\gls{cr} problem} is a trapdoor function where the ``easy part'' is computing exponentiations while the ``hard part'' is computing the $n$-th root of a given number. The \gls{cr} problem is intrinsically a problem only in groups --- which in this case are often derived from fields (\Cref{sec:mathematical_underpinnings:fields}). More formally, given a finite prime field $\mathbb{F}_n$ --- where $n$ is the product of two prime numbers $p$ and $q$, i.e., $n = p \times q$ --- the multiplicative cyclic group of nonzero elements $(\mathbb{F}\textsuperscript{*}\textsubscript{n\textsuperscript{2}}, \times)$, and an element $c \in \mathbb{F}\textsuperscript{*}\textsubscript{n\textsuperscript{2}}$, is it (believed to be) hard to find an element $x \in \mathbb{F}\textsuperscript{*}\textsubscript{n\textsuperscript{2}}$ so that $x^n \equiv c \pmod{n^2}$; computing $x ^ n$ is easy, but recovering $x$ from $c$ --- that is, computing the $n$-th root of $c$ modulo $n$ --- is hard. The word ``residuosity'' refers to the fact that $c$ is the ``residue'' of $x^n \pmod{n^2}$ --- i.e., what remains after applying the modulus $n^2$ to $x^n$ --- while the term ``composite'' refers to the fact that $n$ is composite --- i.e., $n$ is the product of $p$ and $q$. The cyclic group $(\mathbb{F}\textsuperscript{*}\textsubscript{n\textsuperscript{2}}, \times)$ may be built only on top of integers (where the operation is the exponentiation), not on top of elliptic curves (where the operation is point addition) --- see \Cref{sec:mathematical_underpinnings:elliptic_curves}).

Note that not every element $c \in \mathbb{F}\textsuperscript{*}\textsubscript{n\textsuperscript{2}}$ can be written as $x^n \pmod{n^2}$ for a given $x \in \mathbb{F}\textsuperscript{*}\textsubscript{n\textsuperscript{2}}$. In other words, the $n$-th root of a given number may not exist in $\mathbb{F}\textsuperscript{*}\textsubscript{n\textsuperscript{2}}$. Hence, the \gls{cr} problem is often formulated as: given the modulus $n$ and an element $c \in \mathbb{F}\textsuperscript{*}\textsubscript{n\textsuperscript{2}}$, decide whether $c$ is an $n$-residue modulo $n^2$. In other words, it is infeasible to distinguish $n$-th powers from random elements of $(\mathbb{F}\textsuperscript{*}\textsubscript{n\textsuperscript{2}}, \times)$.

\remarkBlock{\gls{cr} and \gls{if}}{The \gls{cr} and \gls{if} (\Cref{sec:trapdoor_functions:integer_factorization}) problems are different albeit strongly connected. In fact, being able to factor $n$ allows for trivially resolving the \gls{cr} problem, and vice versa.}

The \gls{cr} problem is the trapdoor function underlying:
\begin{itemize}
	\item the Paillier cipher (\Cref{sec:asymmetric_cryptography:ciphers:paillier}).
\end{itemize}

\exampleBlock{\gls{cr} example}{Consider the finite prime field $\mathbb{F}_{10}$ and the multiplicative group of nonzero elements $(\mathbb{F}\textsuperscript{*}\textsubscript{100}, \times)$. In this context, it is easy to compute the $10$-residue of $3$ as $3^{10} = 59049 \equiv 49 \pmod{100}$ (hence, $49$ is the $10$-residue). However, given $c = 49$, it is difficult to understand that $49$ is a $10$-residue and $49 \equiv 3 ^{10} \pmod{100}$. In fact, the simplest approach to determine whether an element is a residue consists of trial-and-error computations --- quick for a group with modulus $100$, but infeasible for larger groups. As another example, $c = 2$ is not a $10$-residue, as no element in $\mathbb{F}\textsuperscript{*}\textsubscript{100}$ raised to the 10th power gives $2 \pmod{100}$.}

\warningBlock{The \gls{cr} problem vs. quantum computing}{Quantum computing can solve any \gls{cr} problem using the Shor's algorithm \cite{shor_algorithms_1994} in polynomial time (that is, efficiently) --- see \Cref{sec:TEMPORARY_LABEL:post_quantum_cryptography}.}
